\uuid{G7xK}
\titre{Majorant}
\auteur{bodin ; hanani ; mzari}
\organisation{exo7}
\licence{CC BY-NC-SA 4.0}
\datecreate{2020-01-20}
\niveau{L1}
\module{Analyse}
\chapitre{Propriétés de R}
\sousChapitre{Maximum, minimum, borne supérieure}
\difficulte{2}
\theme{Maximum, majorant, majorant}
\exo7qcmid{lille-1.reels.029}

\contenu{
\begin{question}
\qclassification{120.02}
\qtags{motcle=majorant}
On dit que $M \in \Rr$ est un majorant d'un ensemble $E \subset \Rr$ si pour tout $x\in E$, on a $x \le M$.
\begin{answers}
    \good{Si $E = [3,7]$ alors $8$ est un majorant de $E$.}

    \good{Si $E = [-3,-1[$ alors tout $M \ge -1$ est un majorant de $E$.}

    \bad{Si $E = ]0,+\infty[$ alors tout $M \ge 0$ est un majorant de $E$.}

    \bad{Si $E = [2,3] \cup [5,10]$ alors tout $M \ge 3$ est un majorant de $E$.}
\end{answers}
\begin{explanations}
Tracer les intervalles sur la droite réelle pour mieux comprendre. Les majorants d'un ensemble sont alors tous les réels "à droite" de l'ensemble.
\end{explanations}
\end{question}
}
